% TODO make hw stylesheet
\documentclass[10pt]{article}
\usepackage[letterpaper,top=1in]{geometry}
\usepackage[dvipsnames,svgnames]{xcolor}
\usepackage[shortlabels]{enumitem}
\usepackage[framemethod=TikZ]{mdframed}
\usepackage{amsmath,amssymb,amsthm}
\usepackage[colorlinks]{hyperref}
\usepackage{multicol}
\usepackage{tabularx}

\hypersetup{
	colorlinks=true,
	linkcolor=blue,
	filecolor=magenta,
	urlcolor=magenta,
	pdftitle={MAT 528 HW}
}

\usepackage{mathtools}
\usepackage{thmtools}
\usepackage[textsize=footnotesize]{todonotes}
\usepackage{titling}

\reversemarginpar
\mdfdefinestyle{mdbluebox}{roundcorner=10pt,innerbottommargin=9pt,
	linecolor=blue,backgroundcolor=TealBlue!5,}
\declaretheoremstyle[headfont=\sffamily\bfseries\color{MidnightBlue},
mdframed={style=mdbluebox},]{thmbluebox}

\mdfdefinestyle{mdbloobox}{frametitlefont=\bfseries,innerbottommargin=8pt,
	nobreak=true,backgroundcolor=TealBlue!5,linecolor=blue,}
\declaretheoremstyle[headfont=\bfseries\color{MidnightBlue},
mdframed={style=mdbloobox},headpunct={\\[3pt]},postheadspace=0pt,]{thmbloobox}

\mdfdefinestyle{mdredbox}{frametitlefont=\bfseries,innerbottommargin=8pt,
	nobreak=true,backgroundcolor=Salmon!5,linecolor=RawSienna,}
\declaretheoremstyle[headfont=\bfseries\color{RawSienna},
mdframed={style=mdredbox},headpunct={\\[3pt]},postheadspace=0pt,]{thmredbox}

\mdfdefinestyle{mdgreenbox}{linecolor=ForestGreen,backgroundcolor=ForestGreen!5,
	linewidth=2pt,rightline=false,leftline=true,topline=false,bottomline=false,}
\declaretheoremstyle[headfont=\bfseries\sffamily\color{ForestGreen!70!black},
mdframed={style=mdgreenbox},headpunct={ --- },]{thmgreenbox}

\mdfdefinestyle{mdorangebox}{linecolor=RedOrange,backgroundcolor=RedOrange!5,
	linewidth=2pt,rightline=false,leftline=true,topline=false,bottomline=false,}
\declaretheoremstyle[headfont=\bfseries\sffamily\color{RedOrange!70!black},
mdframed={style=mdorangebox},headpunct={ --- },]{thmorangebox}

\mdfdefinestyle{mdblackbox}{linecolor=black,backgroundcolor=RedViolet!5!gray!5,
	linewidth=3pt,rightline=false,leftline=true,topline=false,bottomline=false,}
\declaretheoremstyle[mdframed={style=mdblackbox}]{thmblackbox}

\declaretheoremstyle[spaceabove=6pt,spacebelow=6pt]{basehead}

\declaretheorem[style=basehead,name=Problem]{problem}
\declaretheorem[style=thmredbox,name=Problem,numbered=no]{problem*}
\declaretheorem[style=thmbloobox,name=Setup]{setup} % for config geo
\declaretheorem[style=thmredbox,name=Example,sibling=problem]{example}
\declaretheorem[style=thmredbox,name=$\bigstar$ Problem,sibling=problem]{niceprob}
\declaretheorem[style=thmbluebox,name=Theorem,sibling=problem]{theorem}
\declaretheorem[style=thmbluebox,name=Corollary,sibling=theorem]{corollary}
\declaretheorem[style=thmbluebox,name=Proposition,sibling=theorem]{prop}
\declaretheorem[style=thmbluebox,name=Lemma,numberwithin=problem]{lemma}
\declaretheorem[style=thmbluebox,name=Theorem,numbered=no]{theorem*}
\declaretheorem[style=thmbluebox,name=Lemma,numbered=no]{lemma*}
\declaretheorem[style=thmgreenbox,name=Claim,numberwithin=problem]{claim}
\declaretheorem[style=thmgreenbox,name=Claim,numbered=no]{claim*}
\declaretheorem[style=thmblackbox,name=Remark,numberwithin=problem]{remark}
\declaretheorem[style=thmblackbox,name=Remark,numbered=no]{remark*}
\declaretheorem[style=thmgreenbox,name=Definition,sibling=theorem]{definition}
\declaretheorem[style=thmgreenbox,name=Definition,numbered=no]{definition*}

\providecommand{\ol}{\overline}
\providecommand{\ceil}[1]{\left\lceil #1 \right\rceil}
\providecommand{\floor}[1]{\left\lfloor #1 \right\rfloor}
\newcommand{\dd}{\mathrm{d}}
\providecommand{\ul}{\underline}
\providecommand{\eps}{\varepsilon}
\providecommand{\half}{\frac{1}{2}}
\providecommand{\CC}{\mathbb C}
\providecommand{\FF}{\mathbb F}
\providecommand{\NN}{\mathbb N}
\providecommand{\QQ}{\mathbb Q}
\providecommand{\RR}{\mathbb R}
\providecommand{\ZZ}{\mathbb Z}
\providecommand{\dg}{^\circ}
\providecommand{\ii}{\item}
\providecommand{\setdiff}{\smallsetminus}
\providecommand{\alert}[1]{{\sffamily\textbf{\textcolor{blue}{#1}}}}
\providecommand{\inv}{^{-1}}
\providecommand{\mb}[1]{\mathbf{#1}}

\newenvironment{soln}{\begin{proof}[Solution]}{\end{proof}}

\providecommand{\pow}{\operatorname{Pow}}
\providecommand{\vocab}[1]{{\textbf{\textcolor{ForestGreen}{#1}}}}
\providecommand{\lcm}{\operatorname{lcm}}
\providecommand{\abs}[1]{\left|#1\right|}

\usepackage{mathrsfs}

% place title at top right
\setlength{\droptitle}{-8ex}
\pretitle{\begin{flushright}\Large\bfseries}
\posttitle{\par\end{flushright}}
\preauthor{\begin{flushright}}
\postauthor{\end{flushright}}
\predate{\begin{flushright}}
\postdate{\end{flushright}}

\title{MAT 528 HW08}
\author{\large Aditya Pahuja}
\date{\large Due 10/30}
\begin{document}
\maketitle

\begin{problem}[Munkres 30.4]
    Show that every compact metrizable space $X$ has a countable basis.
\end{problem}

\begin{soln}
    Let $d$ be a metric on $X$ that induces the topology of $X$.
    For each $n\in\NN$, let $\mathcal U_n$ be the collection of
    open balls in $X$ (with respect to $d$) with radius $1/n$.
    This is an open cover of $X$,
    so since $X$ is compact it has a finite subcover $\mathcal B_n$.
    Let $\mathcal B = \bigcup_{n=1}^\infty\mathcal B_n$;
    this is a union of countably many finite sets, so it is countable.

    I claim that this collection is a basis for the topology on $X$.
    For any $x\in X$ and open $U\subseteq X$ containing $x$,
    let $B_r(x)$ be an open ball centered at $x$ with radius $r$
    contained in $U$. Then, pick positive integer $m$
    such that $\frac 1m < \frac r{2}$. There exists a ball
    $B_{1/m}(x')\in\mathcal B_m$ containing $x$ because $\mathcal B_m$
    is an open cover of $X$, and this ball is a subset of $B_r(x)$ because
    if $y\in B_{1/m}(x')$, then
    \begin{equation*}
        d(x,y) \le d(x,x') + d(x',y) < \frac 1m + \frac 1m
	< \frac r2 + \frac r2 = r.
    \end{equation*}
    Thus given $x\in X$ and $U\subseteq X$ open containing $x$,
    there is an element $B\in \mathcal B$ such that
    $x\in B\subseteq U$, which is what it means
    for $\mathcal B$ to be a basis for the topology on $X$.
\end{soln}

\begin{problem}[Munkres 30.9]
    Let $A$ be a closed subspace of $X$. Show that if $X$ is Lindel\"of,
    then $A$ is Lindel\"of. Show by example that if $X$ has
    a countable dense subset, $A$ need not have a countable dense subset.
\end{problem}

\begin{soln}
    Let $\mathcal U$ be an open cover of $A$.
    Each element $U\in \mathcal U$ is the intersection of an open set $U'$
    with $A$; let $\mathcal U'$ be the collection of these $U'$.
    Then, $\mathcal U' \cup \{X\setminus A\}$ is an open cover of $X$,
    so it has a countable subcover $\mathcal V$.
    If $X\setminus A$ is in this subcover, remove it;
    the resulting collection of open subsets of $X$
    is a countable subset of $\mathcal U'$, and since $X\setminus A$
    is disjoint from $A$, $\mathcal V\setminus\{X\setminus A\}$
    is an open cover of $A$. Thus $\{V\cap A : V\in\mathcal V\}$
    is a countable subcover of $\mathcal U$, proving that $A$ is Lindel\"of.

    $\RR_\ell^2$ is separable with countable dense subset $\QQ^2$,
    since any basis element $[a,b)\times[c,d)$ contains a point
    with rational coordinates. On the other hand, the subset
    $L = \{(a,-a) : a\in\RR\}$ is not separable because it is uncountable
    and inherits the discrete topology from $\RR_\ell^2$,
    as shown in Munkres's Example 30.4.
\end{soln}

\begin{problem}[Munkres 31.8]
    Let $G$ be a compact topological group;
    let $X$ be a topological space;
    let $\alpha$ be an action of $G$ on $X$.
    If $X$ is Hausdorff, regular, normal, locally compact, or second-countable,
    then so is $X/G$.
\end{problem}

\begin{soln}
    We show that the projection map $\pi\colon X\to X/G$
    is closed, surjective, and perfect, so that by Munkres's Exercises 31.6 and 31.7,
    the problem statement holds.

    Since $X\to X/G$ is a quotient map, it is surjective.
    Let $A\subseteq X$ be closed. Then $\pi(A)$ is closed in $X/G$
    if and only if $\pi\inv(\pi(A))$, which is the union of
    the orbits of all elements in $A$, is closed in $X$.
    We can write this set equivalently as $GA\coloneq\{g\cdot x : g\in G, x\in A\}$,
    which is the image $\alpha(G\times A)$.
    If $y\notin GA$ then $gy\notin A$ for all $g\in G$.
    Since $X\setminus A$ is open, for each $g$ there is a neighborhood
    $U_g\times V_g\subseteq G\times X$ of $(g,y)$ such that
    $\alpha(U_g\times V_g)\cap A = \varnothing$.
    The $U_g$ cover $G$, so we can choose a finite subcover
    $U_{g_1}$, \dots, $U_{g_n}$ because $G$ is compact.
    Then $V = V_{g_1}\cap\cdots\cap V_{g_n}$ is a neighborhood of $y$
    and $\alpha(G\times V)$ does not intersect $A$,
    so $V$ does not intersect $GA$ (otherwise, if $z\in V\cap GA$,
    then $g_0z \in A$ for some $g_0\in G$, so $\alpha(G\times\{z\})$ intersects $A$).
    This proves that $GA$ is closed, so $\pi$ is a closed map.

    Finally, $\pi$ is perfect because the orbit of any $x\in X$
    is $Gx\coloneq\{g\times x : g\in G\}$, which is the continuous image
    of the compact set $G$, so the preimage of any point $[Gx]\in X/G$,
    which is $\pi\inv([Gx]) = Gx$, is compact.
\end{soln}

\begin{problem}[Munkres 32.4]
    Show that every regular Lindel\"of space is normal.
\end{problem}

\begin{soln}
    Let $X$ be regular Lindel\"of, and let $A,B\subseteq X$ be
    disjoint closed subsets of $X$.
    Each point $x$ of $A$ has a neighborhood $U$ not intersecting $B$.
    Using regularity, choose a neighborhood $V$ of $x$ whose closure lies in $U$.
    The collection of all such $V$ (as $x$ varies over $A$)
    is an open cover of $A$, so it has a countable subcover $\{U_n : n\in\NN\}$
    because closed subsets of Lindel\"of spaces are Lindel\"of
    (by Munkres's exercise 30.9).
    By a similar procedure, choose a countable collection
    of open sets $\{V_n : n\in\NN\}$ covering $B$
    such that $\ol{V_n}$ is disjoint from $A$ for each $n$.

    Given $n$, define
    \begin{equation*}
	U'_n = U_n\setminus\bigcup_{i=1}^n\ol{V_i},\quad
	V'_n = V_n\setminus\bigcup_{i=1}^n\ol{U_i}.
    \end{equation*}
    Each set $U'_n$ is open because $U_n$ is being intersected with
    the complement of the closed set $\bigcup_{i=1}^n\ol{V_i}$,
    and $U'_n$ covers $A$ because $\ol{V_i}$ is disjoint from $A$ for all $i$.
    Similarly, $\{V'_n : n\in\NN\}$ is an open cover of $B$.
    Finally, by construction, $\bigcup_{n=1}^\infty U'_n$ and
    $\bigcup_{n=1}^\infty V_n$ are disjoint
    since $U'_n$ is disjoint from $V_j\supseteq V'_j$ for all $j\le n$
    and and $V'_n$ is disjoint from $U_j\supseteq U'_j$ for all $j\le n$,
    so we have separated $A$ and $B$ by the open sets $\bigcup_{n=1}^\infty U'_n$
    and $\bigcup_{n=1}^\infty V'_n$.
\end{soln}










\end{document}

